\documentclass{article}
\usepackage[margin=1in]{geometry}
\usepackage{amsmath, amssymb, amsthm}
\usepackage{enumitem}
\usepackage{mathtools}
\usepackage{cancel}

% colored links
\usepackage{hyperref}
\hypersetup{
    colorlinks=true,
    linkcolor=blue,
    filecolor=magenta,      
    urlcolor=black!20!BurntOrange,
    }

% Inputting Python code
\usepackage[dvipsnames]{xcolor}
\definecolor{textblue}{rgb}{.2,.2,.7}
\definecolor{textred}{rgb}{0.54,0,0}
\definecolor{textgreen}{rgb}{0,0.43,0}
\usepackage{upquote}
\usepackage{listings}
\lstset{
    language=Python, 
    tabsize=4,
    basicstyle={\ttfamily},
    keywordstyle=\color{textblue},
    commentstyle=\color{textgreen},
    stringstyle=\color{textred},
    frame=none,
    columns=fullflexible,
    keepspaces=true,
    showstringspaces=false,
    xleftmargin=-15mm, % manual adjustment, figure out permanent solution
}

% Drawing figures
\usepackage{tikz}
\usetikzlibrary{calc, shapes.symbols}

% Colored Boxes
\usepackage{tcolorbox}
\tcbuselibrary{skins,hooks}
\usetikzlibrary{shadows}
\usepackage{lipsum}

%Images
\usepackage{graphicx}
    \usepackage{subcaption}
    \usepackage{float}

%Formatting and Spacing
\setitemize[1]{noitemsep, parsep = 5pt, topsep = 5pt}
\setenumerate[1]{label = (\alph*), parsep = 1pt, topsep = 5pt}
\setlength\parindent{0pt}
\linespread{1.1}

%Easy Numbering in case we add or remove questions
\newcounter{counter}
\newcommand{\Exercise}{Exercise \stepcounter{counter}\thecounter}

\newcommand{\Cov}{\text{Cov}}
\newcommand{\trace}{\text{trace}}
\newcommand{\Normal}{\text{Normal}}
\newcommand{\var}[1]{\operatorname{var}\left(#1\right)}
\newcommand{\E}[1]{\operatorname{E}\left[#1\right]}
% \renewcommand{\Pr}[1]{\operatorname{P}\left(#1\right)}

\DeclareMathOperator*{\minimize}{minimize}
\DeclareMathOperator*{\maximize}{maximize}
\DeclareMathOperator*{\argmin}{argmin}
\DeclareMathOperator*{\argmax}{argmax}
\DeclarePairedDelimiter\abs{\lvert}{\rvert}%
\DeclarePairedDelimiter\norm{\lVert}{\rVert}%
\DeclarePairedDelimiterX{\inp}[2]{\langle}{\rangle}{#1, #2}

%Custom Title Fields
\newcommand{\hmwkTitle}{Homework 1 Solutions}
\newcommand{\hmwkDueDate}{Jan.\ 24, 2024}
\newcommand{\hmwkDueTime}{11:59pm EST}
\newcommand{\hmwkClass}{Advanced Algorithms}
\newcommand{\hmwkClassInstructor}{Professor Dana Randall}
\newcommand{\hmwkSection}{Spring 2024}
\newcommand{\hmwkAuthorName}{}

%Headers and Footers
\usepackage{fancyhdr}
\usepackage{extramarks}
\pagestyle{fancy} 
\lhead{CS 4540}
\chead{\hmwkClass \ (\hmwkClassInstructor)}
\rhead{\hmwkTitle}
\cfoot{\thepage}
\renewcommand\headrulewidth{0.4pt}
\renewcommand\footrulewidth{0.4pt}

\title{
    \vspace{2in}
    \textbf{\hmwkClass:\ \hmwkTitle}\\
    \normalsize\vspace{0.1in}\small{Due\ on\ \hmwkDueDate\ at \hmwkDueTime}\\
    \vspace{0.1in}\large{\textit{\hmwkClassInstructor\ \hmwkSection}} \\
    \vspace{0.8in}
    \begin{minipage}[t]{0.4\columnwidth}
        {\footnotesize\textbf{As stated in the syllabus, unauthorized use of previous semester course materials is strictly prohibited in this course.
        }}
    \vspace{0.5in}
    \end{minipage}
    \vspace{0.8in}
    \author{\textbf{\hmwkAuthorName}}
    \date{}
}

\begin{document}

\maketitle
\pagebreak


\subsection*{\Exercise}
    \textit{(KT Chapter 1, Exercise 1)} True of false?  (If true, give a brief explanation, if false, give a counterexample): 
    
    In every instance of the Stable Matching Problem, there is a stable matching containing a pair $(m, w)$ such that $m$ is ranked first on the preference list of $w$ and $w$ is ranked first on the preference list of $m$.

\subsection*{Solution}
    This is false. Here is an example with 4 people: $m,m'$ and
    $w,w'$. Say $m$'s list is $[w,w']$, $m'$'s list is $[w',w]$. However, $w$'s list is $[m',m]$, and $w'$'s list is $[m,m']$.

    \vspace{2mm}
    There are two possible pairings, and both are stable. Either both men get their first choice and both women get their second choice, or vice-versa. In neither pairing is there a marriage where both partners rank the other first on their preference lists. 

\subsection*{\Exercise}
    \textit{(KT Chapter 1, Exercise 2)} True or false?
    
    Consider an instance of the Stable Matching Problem in which there exists a man $m$ and a woman $w$ such that $m$ is ranked first on the preference list of $w$ and $w$ is ranked first on the preference list of $m$.  Then in every stable matching $S$ for this instance, the pair $(m, w)$ belongs to $S$.

\subsection*{Solution}
    This is true. Assume otherwise - say there is a stable matching $M$ where $w$ is not matched to $m$, and $m$ is not matched to $w$. Since each is first on the other's list, and they are currently matched to someone who is not first on their list, $(m,w)$ is a rogue pair. Thus $M$ isn't stable - contradiction!.

\subsection*{\Exercise}
    \textit{(KT Chapter 1, Exercise 3)} There are many other settings in which we can ask questions related to some type of ``stability'' principle.  Here's one, involving competition between two enterprises.  

    \vspace{3mm}
    Suppose we have two television networks whom we'll call $A$ and $B$.  There are $n$ prime-time
    programming slots, and each network has $n$ TV shows.  Each network wants to devise a schedule -- an assignment of each show to a distinct slot -- so as to attract as much market share as possible.

    \vspace{3mm}
    Each show has a fixed rating, which is based on the number of people who watched it last year;  we'll assume that no two shows have exactly the same rating.  A network {\it wins} a given time slot if the show that it schedules for that time slot has a larger rating than the other network schedules for that time slot.  The goal is to win as many time slots as possible.

    \vspace{3mm}
    We'll say a pair of schedules $(S, T)$ is {\it stable} if neither network can unilaterally change its own schedule and win more time slots.  

    \vspace{3mm}
    For every set of TV shows and ratings, is there always a stable pair of schedules?  Resolve this question by doing one of two things: 

    \begin{enumerate}
        \item give an algorithm that, for any TV shows and associated ratings, produces a pair of stable schedules; or
        \item give an example of TV shows and ratings for which there is no stable pair of schedules.  
    \end{enumerate}

\subsection*{Solution}
    Unlike stable matching, there isn't always a stable schedule assignment. Say one network A has shows with ratings $1,3$, and the other network B has shows with ratings $2,4$. If $B$ had it's choice, it would put $2$ at the same slot as $1$, and $4$ at the same slot as $3$, so that it would win both time slots. But this isn't stable, as $A$ would swap it's schedule so that $1$ is against $4$, and $3$ is against $2$, and each network would win one time slot. This isn't stable either, as with this starting configuration, $B$ would swap it's shows to win both time slots.

\subsection*{\Exercise}
    Give an input to the Stable Marriage Problem with $n$ men and $n$ women that has a unique stable pairing.   Now give an input with $n$ men and $n$ women that has an exponential number (in $n$) of stable pairings.

\subsection*{Solution}
    \begin{itemize}
      \item One situation was presented in class: every man and every woman have
        the same preference list. $[m_1,m_2,\ldots]$ for each woman, and
        $[w_1,w_2,\ldots]$ for each man. Then, the only stable matching is
        $(m_1,w_1), (m_2,w_2),\ldots$. \underline{Proof}: Assume there is another matching $M$.
        Consider the smallest $i$ so that $m_i$ isn't matched with $w_i$.
        Then each of these is matched with someone lower on their preference list, and thus $(m_i,w_i)$ is
        a rogue pair - contradiction.
    
        Another situation with a unique stable matching is the following:
        say for all i, each person $m_i$ has
        $w_i$ as their ideal mate, and each $w_i$ has $m_i$ as their ideal mate.
        Then the matching $(m_1,w_1), (m_2,w_2),\ldots$ is stable, since every
        person is paired with their ideal mate. Considering problem $2$, this is
        the only possible stable matching. Note that it doesn't matter what
        all the other preferences are.
    
      \item There are many possible ways to get an exponential number of stable
        matchings, but here is one.
    
        Say $n$ is even, divide all the men and all the women into groups of 2.
        Label the men $m_1,m_1', m_2,m_2',\ldots m_{n/2},m_{n/2}'$, and do the
        same for the women.
    
        For each pair, we use a setup like the written solution from problem 1:
        Let each $m_i$'s list be $[w_i,w_i',\ldots]$, $m_i'$'s list be
        $[w_i',w_i,\ldots]$, and let each
        $w_i$'s list be $[m_i',m_i,\ldots]$, and $w_i'$'s list be
        $[m_i,m_i',\ldots]$. All other preferences after these first two won't
        matter.
    
        For each $i$, say that we either have pairing $[(m_i,w_i), (m_i',w_i')]$,
        or pairing $[(m_i, w_i'),(m_i',w_i)]$. 
        We can see that this is stable because 
        \begin{itemize}
          \item Everyone is paired with either their first or second choice, so
            there are no rogue pairs between some $m_i$ and some $w_j$, $j\neq i$.
          \item No one can move unilaterally within their pair. As in the solution
            to problem 1, this is stable.
        \end{itemize}
    
        We can select either pair for each $i$ without affecting any of the
        others. Thus there are at least $2^{\frac{n}{2}}$ stable matchings.
        (Without much more effort, you can show that these are the only stable
        matchings - but this is unnecessary).
    \end{itemize}

\subsection*{\Exercise}
    Let $[m]$ denote the set $\{0,1,\dots, m-1\}$.  For each of the following families of hash functions, say whether it is 2-universal or not, and specify how many random bits are needed to sample a hash function from the family, and then justify your answer.
    \begin{enumerate}
        \item $H_a = \{h(x_1,x_2)=a_1 x_1 + a_2 x_2 \ ({\rm mod} \ m) \ | \ a_1, a_2 \in [m]\}$, where $m$ is some fixed prime.  (Note that for each $h\in H_a$ we have $h: [m]^2 \rightarrow [m]$, i.e., it maps a pair of integers in $[m]$ to a single integer in $[m]$.)
        \item $H_b = \{h(x_1,x_2) = a_1x_1 + a_2 x_2 \ ({\rm mod} \  m) \} : a_1, a_2 \in [m]\}$, where $m=2^k$ is some fixed power of two.
        \item The set of all functions $f: [m]\rightarrow [m-1]$.
    \end{enumerate}

\subsection*{Solution}
\begin{enumerate}[label=(\alph*)]
  \item This is a 2-universal family.
    We use two important facts from number theory that we also used in class.
    \begin{itemize}
      \item Fact 1: Let $a$ be uniformly chosen from $[m]$. Then for ANY fixed
        $x$, $a + x \pmod m$ is uniformly distributed in $[m]$.
      \item Fact 2: Let $a$ be uniformly chosen from $[m]$, where $m$ is prime.
        Then for any fixed $x \neq 0$, $a\cdot x \pmod m$ is
        uniformly distributed in $[m]$.
    \end{itemize}

    This hash family is 2-universal if $$\forall (x_1,x_2) \neq (y_1,y_2), 
    \left[\Pr[f(x_1,x_2) = f(y_1,y_2)] \leq 1/m\right].$$
    Notice that the $\forall$ is outside of the $\Pr$, so we
    \emph{first} fix $(x_1,x_2)\neq (y_1,y_2)$, \emph{then} we consider the
    Probability over all $a_1,a_2$.

    The event $f(x_1,x_2) = f(y_1,y_2)$ means that
    $a_1x_1 + a_2x_2 = a_1y_1 + a_2y_2 \pmod m$. With some algebra, it means
    $a_1(x_1 - y_1) + a_2(x_2 - y_2) = 0 \pmod m$.
    Since $(x_1,x_2) \neq (y_1,y_2)$, we know at least one of $x_1\neq y_1$ or
    $x_2 \neq y_2$. Without loss of generality, assume that $x_1 \neq y_1$.
    Then, $(x_1 - y_1)$ is non-zero. Thus, $a_1(x_1 - y_1) \pmod m$ is
    distributed uniformly in $[m]$ by Fact (2).
    The other term $a_2(x_2 - y_2) \pmod m$ is something independent of
    $a_1$. No matter what it $a_2(x_2 - y_2)$ happens to be,
    $a_1(x_1 - y_1) + a_2(x_2 - y_2) \pmod m$ is
    uniformly distributed in $[m]$ by Fact (1). Thus the the Probability that
    this is $0$ is $1/m$.

    A hash function in this family is determined by $a_1,a_2$, each of which
    are $\in [p]$. Thus we need $2\lceil\log(p)\rceil$ bits.

  \item This is NOT a 2-universal family.
    Now that $m = 2^k$ is not Prime, we can no longer use Fact (2) in our Proof
    above.
    Consider $(x_1,x_2) = (0,0)$ and $(y_1,y_2) = (2^{k-1},0)$. Now, the event
    $f(x_1,x_2) = f(y_1,y_2)$ can be boiled down with some algebra to
    $a_12^{k-1} = 0 (\mod 2^k)$. If $a_1$ is a multiple of $2$, then $a_12^{k-1}$
    is a multiple of $2^k$. Exactly half of the elements of $[2^k]$ are
    multiples of 2.
    Thus for any $k > 1$, $\Pr[f(0,0) = f(2^{k-1},0)]=\frac{1}{2} > \frac{1}{2^k}$.

  \item This is a 2-universal family.
    In this family, each input $i \in [m]$ has $m-1$ possible choices to map
    to.  We can sample from this set in the following natural way - 
    for each element $i \in [m]$, pick an element in $[m-1]$ uniformly at
    random, independently of our choices for all other $j \in [m]$.

    Therefore, for any fixed $x,y \in [m]$, $f(x)$ and $f(y)$ are
    independent, and $\Pr[f(x) = f(y)] = \frac{1}{m-1}$. Since our target set
    has size $m-1$, this is a 2-universal family.

    We need to choose one element in $[m-1]$ for each element in $[m]$. Thus
    we need $m\lceil\log(m-1)\rceil$ bits.

\end{enumerate}



\end{document}
